\documentclass[11pt]{article}
\usepackage[a4paper,margin=27mm]{geometry}
\usepackage{amsmath,amssymb,amsthm,hyperref}
\newtheorem{theorem}{Theorem}
\newtheorem{lemma}[theorem]{Lemma}
\newtheorem{corollary}[theorem]{Corollary}
\newcommand{\Q}{\mathbb Q}
\newcommand{\Z}{\mathbb Z}
\title{Optimizing the parity filter by a null polynomial}
\author{Research note, 30 September 2026}
\date{}
\begin{document}
\maketitle

\begin{theorem}\label{main}
If rational numbers $x_1,\ldots,x_K$ satisfy
\[
 \frac1K\sum_{i=1}^K x_i^j=\frac1{j+1}\qquad(0\le j\le r),
\]
then
\[
 K\ge 2^{2\lfloor r/2\rfloor}.
\]
The nodes may repeat and need not belong to $[0,1]$.
\end{theorem}

Write $v=v_2$, and recall from
\path{../rational_designs/two_adic_lower_bound.tex} that
\[
 G(z)=\frac z{\log(1+z)}=\sum_{j\ge0}g_jz^j,
 \qquad v(g_j)=-j.
\]
For completeness, $G(-2t)^{-1}=\sum_{j\ge0}2^jt^j/(j+1)$
has integral coefficients and is $1+t$ modulo $2$.
Thus $G(-2t)\equiv(1+t)^{-1}\pmod2$, proving the valuation assertion.
Also $g_j=\int_0^1\binom{x}{j}\,dx$.

\begin{lemma}\label{dilation}
Let $N=2^a$ with $a\ge1$, and put
\[
 I_j(N)=\int_0^1\binom{Nx}{j}\,dx.
\]
Then
\[
 v(I_j(N))\ge-\lfloor j/N\rfloor,
 \qquad v(I_{N\ell}(N))=-\ell\quad(\ell\ge0).
\]
\end{lemma}
\begin{proof}
Coefficientwise integration gives
\[
 \sum_{j\ge0}I_j(N)z^j=G(W(z)),
 \qquad W(z)=(1+z)^N-1.
\]
Every coefficient of $W$ below degree $N$ is even, whereas its leading
coefficient is one. Consider a contribution from $g_kW^k$ to degree $j$,
and let $b$ be the number of its factors having degree less than $N$.
It has valuation at least $b-k$. If $L=\lfloor j/N\rfloor$, at most
$L$ factors can have degree $N$, hence $k-b\le L$. This proves the bound.

For $j=N\ell$, no term with $k<\ell$ occurs. The term with
$k=\ell$ is uniquely $g_\ell z^{N\ell}$ and has valuation $-\ell$.
If $k>\ell$, having $\ell$ factors of degree $N$ already leaves no room
for the other positive-degree factors. Thus $k-b\le\ell-1$, and all
these contributions have valuation at least $1-\ell$. They cannot
cancel the unique term of valuation $-\ell$.
\end{proof}

The classical parity filter is
\[
 H_s(y)=\sum_{j=1}^s(-2)^{j-1}\binom yj.
\]
For $y\in\Z_2$, the binomial expansion of $(-1)^y$, first for
nonnegative integers and then by continuity, gives
\[
 H_s(y)\equiv\boldsymbol1_{y\notin2\Z_2}\pmod{2^s}.
\]
This filter is attributed to Hua in Mit'kin, Lemma~7 and the following
remark; see the historical discussion and primary source in the note
cited above. The modification below is the additional step in this
derivation; its novelty in the literature has not been audited.

\begin{lemma}\label{filter}
For every even $s\ge2$ and every $N=2^a$, $a\ge1$, there is a rational
polynomial $P$ of degree at most $s$ satisfying
\[
 P(y)\equiv\boldsymbol1_{y\notin2\Z_2}\pmod{2^s}
 \quad(y\in\Z_2),
 \qquad \int_0^1P(Nx)\,dx=0.
\]
\end{lemma}
\begin{proof}
Let $J=\int_0^1H_s(Nx)\,dx$ and $L=\lfloor s/N\rfloor$.
By Lemma~\ref{dilation}, the coefficients
$A_j=(-2)^jI_j(N)$ obey
\[
 v(A_j)\ge j-\lfloor j/N\rfloor.
\]
The series converges at $z=-2$. Moreover,
$W(-2)=(-1)^N-1=0$, so
\[
 \sum_{j\ge0}A_j=G(W(-2))=1.
\]
This evaluation is legitimate directly from the coefficient bound:
the contributions to degree $j$ in the formal composition have the same
lower bound, which tends to infinity. Hence regrouping is valid.
Consequently
\[
 J=-\frac12\sum_{j=1}^s A_j
   =\frac12\sum_{j>s}A_j,
 \qquad
 v(J)\ge s-\lfloor(s+1)/N\rfloor=s-L.
\]
The last equality uses that $s$ and $N$ are even, so $N$ cannot divide
$s+1$. Put $j_0=NL\le s$; this includes $j_0=0$ if $N>s$.
Lemma~\ref{dilation} gives $v(I_{j_0}(N))=-L$, so the rational number
\[
 t=-\frac{J}{2^sI_{j_0}(N)}
\]
belongs to $\Z_2$. Now set
\[
 P(y)=H_s(y)+2^st\binom y{j_0}.
\]
The added polynomial takes values in $2^s\Z_2$ on $\Z_2$, and by
construction its integral cancels $J$ exactly.
\end{proof}

\begin{proof}[Proof of Theorem~\ref{main}]
Take $s=2\lfloor r/2\rfloor$; the case $s=0$ is immediate.
If every $x_i$ is $2$-adically integral, exactness for
$\binom{x}{s}$ and $v(g_s)=-s$ imply $2^s\mid K$.
Otherwise choose the least $a\ge1$ for which every
$y_i=2^a x_i$ lies in $\Z_2$. At least one $y_i$ is odd.
For the polynomial in Lemma~\ref{filter}, exactness gives
\[
 \sum_{i=1}^K P(y_i)=K\int_0^1P(2^a x)\,dx=0.
\]
The number of odd $y_i$ is therefore a positive multiple of $2^s$.
It is at most $K$, proving the theorem.
\end{proof}

\begin{corollary}[Equality forces odd denominators]\label{equality}
If $r\ge2$ is even and $K=2^r$, then every node belongs to $\Z_2$.
Equivalently, each node has odd denominator in lowest terms.
\end{corollary}
\begin{proof}
Otherwise normalize by $N=2^a$ as before. The positive unit count is
a multiple of $2^r=K$, so all $y_i=Nx_i$ are odd. Consequently
$z_i=(y_i-1)/2\in\Z_2$.
Put $M=2^{a-1}$. We claim
\[
 v\left(\int_0^1\binom{Mx-1/2}{r}\,dx\right)
 =-r-v(r!).
\]
For $a\ge2$, expansion around $-1/2$ gives, for every $j\ge0$,
\[
 \int_0^1(Mx-1/2)^j\,dx
 =(-1/2)^j\sum_{\ell=0}^j
       \binom j\ell\frac{(-2M)^\ell}{\ell+1}.
\]
Every summand with $\ell>0$ inside the sum is even, since
$a\ell-v(\ell+1)\ge1$. Thus this integral has valuation $-j$.
For $a=1$, the odd moments vanish and the even moments have valuation
$-j$, by symmetry about $1/2$.
As $r$ is even, the monic degree-$r$ numerator of
$\binom ur$ therefore has a unique lowest-valuation integrated term,
namely $u^r$. This proves the claim.
Exactness and the integrality of every $\binom{z_i}{r}$ now force
$v(K)\ge r+v(r!)>r$, contradicting $K=2^r$.
\end{proof}

\begin{corollary}[An odd-degree refinement]\label{odd}
For odd $r\ge3$, let $o$ be the odd part of $r+1$. Then
\[
 K\ge 2^{r-1}+2^{o-1}.
\]
In particular, degree $5$ requires $K\ge20$, and degree $9$ requires
$K\ge272$.
\end{corollary}
\begin{proof}
If all nodes lie in $\Z_2$, then $2^r\mid K$, which is stronger.
Otherwise take the same minimal normalization $N=2^a$ and let $U$ be
the positive number of odd normalized nodes. Assume $K<2^r$.
The degree-$(r-1)$ even filter shows that $U=2^{r-1}$.

Write $J=\int_0^1H_r(Nx)\,dx$. If $N\nmid r+1$, then
$\lfloor(r+1)/N\rfloor=\lfloor r/N\rfloor$; the proof of
Lemma~\ref{filter} works unchanged in degree $r$ and gives
$2^r\mid U$, a contradiction. Thus $N\mid r+1$; put $q=(r+1)/N$.
The first term in the tail representation of $J$ is
$A_{r+1}/2$, with valuation exactly $r-q$ by Lemma~\ref{dilation}.
Every later term has strictly greater valuation, since
$j-\lfloor j/N\rfloor>(r+1)-q$ for $j>r+1$.
Hence $v(J)=r-q$. Exactness and the parity congruence give
\[
 U\equiv KJ\pmod{2^r}.
\]
As $v(U)=r-1$, necessarily $v(K)=q-1$.
We have $q\ge o$ and $q-1\le(r-1)/2<r-1$. Thus $K$ is strictly
larger than $2^{r-1}$, and the next integer above that value with
valuation $q-1$ is $2^{r-1}+2^{q-1}$. This proves the stated bound.
\end{proof}

\begin{corollary}[The exact minimum in degree five]
The minimum number of nodes in a rational equal-weight uniform-interval
quadrature of degree five is exactly $20$.
\end{corollary}
\begin{proof}
Corollary~\ref{odd} gives the lower bound $20$.
For the upper bound, take the following multiset of ten magnitudes:
\[
 (s_1,\ldots,s_{10})=(1,11,2,2,6,6,8,12,12,14),
\]
and the twenty rational nodes
\[
 x_{i,\pm}=\frac{15\pm s_i}{30}\qquad(1\le i\le10).
\]
They all lie strictly between zero and one.
On the centered interval $[-1,1]$, these nodes are $\pm s_i/15$.
Their odd moments vanish by symmetry, while
\[
 \sum_i s_i^2=750,\qquad \sum_i s_i^4=101250
\]
give the second and fourth moments $1/3$ and $1/5$ exactly.
The affine change of variables proves exactness through degree five.
\end{proof}

\paragraph{Example and scope.}
For $s=2,N=2$, one may take $P(y)=4y-3y^2$.
Its values are the parity indicator modulo $4$, and
$\int_0^1P(2x)\,dx=0$.
The theorem concerns rational scalar quadrature; it does not supply an
exponential lower bound for arbitrary permutation-fair dice.

\end{document}
